\section{Verification from graph derivatives}
\label{sec:schur-seam-verifier}

The main theorem can be applied by checking its assumptions directly, as in
the Hamilton example.  In more complicated problems, an explicit formula
for \(J_{\lambda F}\) may be unavailable.  A pair of smooth parametrizations
of the operator graph can then provide a useful alternative.  We describe
this verification route here; its full statement and proof are in
Appendix~\ref{sec:signed-schur-details}.

Suppose the relevant graph consists of two parametrized branches
\[
 g_\sigma(a,t)=(u_\sigma(a,t),v_\sigma(a,t)),
 \qquad \sigma\in\{+,-\},\quad t\ge0,
\]
which agree at \(t=0\).  Their proximal inputs are
\[
 X_\sigma(a,t)=u_\sigma(a,t)+\lambda v_\sigma(a,t).
\]
The parameter \(a\) runs along the seam, while \(t\) measures displacement
from it on each branch.  The method first solves the tangential components
of \(X_\sigma(a,t)=x\) for \(a\).  The remaining scalar input coordinate
is denoted by \(n_\sigma(z,t)\), where \(z\) is the fixed tangential input.
The key hypothesis requires the two branches to leave their common input
seam in opposite directions:
\[
 \partial_t n_+(z,t)\ge p c_+t^{p-1},\qquad
 -\partial_t n_-(z,t)\ge p c_-t^{p-1},
 \qquad p\ge1,\quad c_+,c_->0.
\]
These inequalities are checked uniformly over the parameter domain.  The
derivative remaining after tangential elimination is a Schur complement;
the two signs distinguish the sides of the seam.

Under the tangential invertibility and derivative bounds of
Theorem~\ref{thm:signed-schur-growth}, together with its complete-fiber
condition, this calculation gives an open input neighborhood \(U\) with
the following properties.  Every \(x\in U\) has exactly one localized
proximal output \(Jx\), and comparisons across the two branches satisfy
\begin{equation}
 \|Jx-Jx'\|^2+
 \|(x-Jx)-(x'-Jx')\|^2
 \le\frac12\left(\|x-x'\|^2+
                     L_D^2\|x-x'\|^{2/p}\right)
 \label{eq:schur-main-J-estimate}
\end{equation}
whenever \(x,x'\) belong to a covered input block of diameter at most
\(D\).  Thus the produced RL exponent is \(\gamma=1/p\).
The constant is explicit: for \(p>1\), put \(p'=p/(p-1)\); then
\[
 L_D=K_zD^{1-1/p}+(1+B^2)^{1/(2p)}
 \left[(K_+c_+^{-1/p})^{p'}+
       (K_-c_-^{-1/p})^{p'}\right]^{1/p'}.
\]
Here \(B,K_z,K_+,K_-\) are the uniform derivative bounds specified in
\eqref{eq:schur-normal-tangent}--\eqref{eq:schur-output-normal}.
Theorem~\ref{thm:signed-schur-growth} also gives the size of \(U\), and
Proposition~\ref{prop:schur-matched-jet} supplies an alternative constant
when the two branch derivatives match to higher order; the leading
H\"older coefficient improves for \(p>1\) with two nonzero branch bounds.

Completeness of the represented graph is part of the verification: every
graph point with an input in \(U\) must be accounted for.  Opposite signed
growth then rules out two distinct outputs at the same input, including
outputs from different branches.  The theorem concerns the localized
resolvent unless completeness is established for the unrestricted graph.

This procedure supplies coverage, uniqueness, and the J-side geometry.
The residual error bound is obtained separately.  To invoke the convergence
theorem one must also verify the nearest-zero anchors in the same graph,
the true-residual evaluation window, scalar compatibility, and the strict
initial length budget.  The proof uses tangential
inversion and one-dimensional growth; its detailed hypotheses and constants
are retained in the appendix for applications in which branch data are
easier to obtain than the resolvent itself.
